\section{\texorpdfstring{Translating the definition into a Lean \texttt{structure}}{Translating the definition into a Lean structure}}\label{sec:06-groups:02-translating:1}

The construction proceeds piece by piece rather than all at once.

\textbf{Step 1 --- just the data}, no axioms yet:

\begin{lstlisting}
structure GroupData (G : Type) where
  op : G → G → G
  id : G
  inv : G → G
\end{lstlisting}

This says: to build a \texttt{GroupData G}, supply an operation, an identity element, and an inverse function. Lean does not check any axioms yet. A \texttt{GroupData} could still be complete nonsense (for example, an \texttt{op} that ignores its arguments). We fix that next.

\begin{mathreading}

\texttt{GroupData G} is the \emph{signature} of a group on the carrier \(G\): the raw data \((\,\cdot : G\times G \to G,\ e \in G,\ (-)^{-1}
: G \to G\,)\) with no laws imposed. As a set it is the product \[
\mathrm{GroupData}(G) = (G^{G\times G}) \times G \times (G^{G}),
\] the collection of all ``magma-with-distinguished-element-and-unary-op'' structures on \(G\). It is a \emph{magma} signature, not yet a group. The axioms below pick out the genuine group structures as a subset.

\end{mathreading}

\textbf{Step 2 --- add the axioms as extra fields.} In Lean, a \emph{proof} can be a field of a structure, just like data can. We add one field per axiom, each of type ``a proposition that must hold for every element'':

\begin{lstlisting}
structure Group (G : Type) where
  op : G → G → G
  id : G
  inv : G → G
  assoc : ∀ a b c : G, op (op a b) c = op a (op b c)
  id_left : ∀ a : G, op id a = a
  id_right : ∀ a : G, op a id = a
  inv_left : ∀ a : G, op (inv a) a = id
  inv_right : ∀ a : G, op a (inv a) = id
\end{lstlisting}

Reading this line by line:

\begin{itemize}
\tightlist
\item
  \texttt{op : G → G → G} --- the group operation, curried (Chapter 2).
\item
  \texttt{id : G} --- a single, fixed element of \texttt{G}, playing the role of \(e\).
\item
  \texttt{inv : G → G} --- a function that assigns to each element its inverse.
\item
  \texttt{assoc : ∀ a b c : G, op (op a b) c = op a (op b c)} --- a \emph{proof obligation}: whoever constructs a \texttt{Group G} must supply a term that proves this statement holds for every \texttt{a b c}.
\item
  The remaining four fields are the identity and inverse axioms. They are split into left and right versions because commutativity has not been assumed.
\end{itemize}

This is the general recipe used throughout the book: \textbf{a mathematical structure is data plus proofs, bundled together}, and Lean's \texttt{structure} mechanism is a direct, literal translation of that idea.

\begin{mathreading}

\texttt{Group G} is the type of \emph{group structures on the fixed carrier \(G\)}: a dependent tuple \[
\Big(\,\cdot,\ e,\ (-)^{-1},\ \underbrace{\alpha}_{\text{assoc}},\
\underbrace{\lambda_\ell,\lambda_r}_{\text{identity}},\
\underbrace{\iota_\ell,\iota_r}_{\text{inverse}}\,\Big),
\] where the last five components are \emph{proofs} of the group axioms, that is, elements of the propositions \[
\forall a,b,c,\ (a\cdot b)\cdot c = a\cdot(b\cdot c);\quad
\forall a,\ e\cdot a = a;\ \ldots;\quad \forall a,\ a\cdot a^{-1}=e.
\] These fields are propositions (proof-irrelevant), so \texttt{Group G} is exactly the \emph{subset} of \texttt{GroupData G} cut out by the group axioms: the \(\Sigma\)-type \(\sum_{d : \mathrm{GroupData}(G)} \mathrm{Axioms}(d)\). This is the same \(\Sigma\)-type already encountered in Chapter 3 as ``a witness together with a proof about it.'' Here the witness is a whole bundle of group data \texttt{d : GroupData G} rather than a single element, and the proof is the combination of the five axioms about it. The shape is identical: it is exactly the pairing a \texttt{structure} performs when it bundles data fields with proof fields, just as \texttt{Group} itself does above. This matches the mathematician's way of saying ``a group is a set with operations \emph{such that} the axioms hold'': the ``such that'' is a genuine \hyperref[sec:01-basics:04-terminology:1]{subobject} of the space of raw data.

\end{mathreading}

\begin{quote}
Read more: Mathlib's actual \texttt{Group} (\texttt{Mathlib.Algebra.Group.Defs}) is a \texttt{class}, not this book's plain \texttt{structure}. It inherits from a chain of smaller classes (\texttt{Mul}, \texttt{One}, \texttt{Inv}, \texttt{Monoid}, \ldots) instead of listing all axioms in one place. See \hyperref[sec:13-next-steps:02-moving-to-mathlib:1]{Chapter 13} for the bridge between the two styles, and \hyperref[sec:05-rigor-check:01-structure-vs-class:1]{Chapter 5 §1} for why this book delays that mechanism.
\end{quote}
